% Part: propositional-logic
% Chapter: syntax-and-semantics
% Section: introduction

\documentclass[../../../include/open-logic-section]{subfiles}

\begin{document}

\olfileid{pl}{syn}{int}

\olsection{پېژندنه}

بياني منطق له هغو !!{formula}s سره کار لري چې له
!!{propositional variable}s څخه د $\lnot$، $\land$، $\lor$، $\lif$ او
$\liff$ بياني نښلوونکو په کارولو جوړېږي۔ د تصور له مخې،
!!a{propositional variable}~$p$ د داسې يوې جملې يا بيان ځاے نيسي
چې رښتونے يا دروغژن وي۔ هر کله چې په !!a{formula} کښې د
!!{propositional variable} ``د رښتياوالي قيمت'' وټاکل شي، نو د
بياني نښلوونکو په کارولو ترې د جوړو شويو ټولو !!{formula}s د
رښتياوالي قيمت هم ټاکل کېږي۔ وايو چې بياني منطق \emph{د رښتياوالي
  تابعوي} دے، ځکه چې معنٰی پوهنه ئې د رښتياوالي د قيمتونو د تابعو
له مخې ورکول کېږي۔ په تېره بيا، په بياني منطق کښې دا نۀ څېړو چې
رښتياوالے او دروغوالے نور څنګه ټاکل کېږي؛ مثلاً دا چې يو څه د يوازې
اتفاقي رښتياوالي پر ځاے بالضروره رښتيا دي، يا څوک پرې پوهېږي چې
رښتيا دي، يا اوس رښتيا دي نۀ دا چې پخوا رښتيا وو يا وروسته به رښتيا
وي۔ موږ يوازې د رښتيا ($\True$) او دروغو ($\False$) دوه قيمتونه په
پام کښې نيسو، نو دا امکان له بحثه وباسو چې يو بيان دې نۀ رښتونے وي
نۀ دروغژن، يا دې نيم رښتونے وي۔ موږ يوازې هغو نښلوونکو ته هم پام
کوو چې ترې د جوړ شوي !!a{formula} د رښتياوالي قيمت د هغې د برخو د
رښتياوالي په قيمتونو پوره ټاکل کېږي، نۀ مثلاً د هغې په معنٰی۔ په
تېره بيا، دا لانجمنه ده چې په انګرېزۍ کښې د شرطيو رښتياوالي قيمت په
دې معنٰی تابعوي دے که نۀ۔ مادي شرطي~$\lif$ تابعوي دے؛ نور منطقونه
له هغو شرطيو سره کار کوي چې د رښتياوالي تابعوي نۀ وي۔

د رښتياوالي د تابعوي بياني منطق د نظريې او مابعدمنطق د ودې دپاره
بايد لومړے د هغه د عبارتونو نحو او معنٰی پوهنه تعريف کړو۔ موږ به له
!!{propositional variable}s څخه د نښلوونکو په کارولو د !!{formula}s
د جوړولو يوه لار بيان کړو۔ نور تعريفونه هم ممکن دي۔ نور نظامونه به
بېلې نښې غوره کړي، د بنسټيزو نښلوونکو بېل سټونه به وټاکي او قوسونه
به په بېله توګه وکاروي، يا به ئې هېڅ و نۀ کاروي، لکه په تش په نامه
پولنډۍ نښه‌اېښودنه کښې۔ خو په ټولو لارو کښې يو ګډ شے دا دے چې د
جوړېدو قاعدې د !!{formula}s سټ \emph{په استقرايي ډول} تعريفوي۔ که
کار په سمه توګه وشي، هر عبارت د جوړېدو د قاعدو له مخې په اصل کښې
يوازې په يوه لار ترلاسه کېدے شي۔ هغه استقرايي تعريف چې
\emph{په يکتا ډول لوستل کېدونکي} عبارتونه ورکوي، موږ ته اجازه راکوي
چې په هماغه طريقه، يعنې په استقرايي تعريف، دغو عبارتونو ته معنٰی
ورکړو۔

عبارتونو ته معنٰی ورکول د معنٰی پوهنې ډګر دے۔ د بياني منطق په
معنٰی پوهنه کښې مرکزي مفهوم په !!a{valuation} کښې صدق دے۔
!!^a{valuation}~$\pAssign{v}$ !!{propositional variable}s ته د
رښتياوالي $\True$ او $\False$ قيمتونه ورکوي۔ هر !!{valuation} د هر
!!{formula}~$!A$ دپاره د رښتياوالي قيمت $\pValue{v}(!A)$ ټاکي۔
!!^a{formula} په !!a{valuation}~$\pAssign{v}$ کښې هله او يوازې هله
صادق دے چې $\pValue{v}(!A) = \True$ وي؛ دا په $\pSat{v}{!A}$ ليکو۔
دا اړيکه د~$!A$ پر جوړښت د استقرا له مخې هم تعريفېدے شي: د منطقي
نښلوونکو د رښتياوالي تابعې کاروو، څو مثلاً د $!A \land !B$ صدق د
$!A$ او~$!B$ د صدق يا نۀ صدق له مخې تعريف کړو۔

د جملو د صدق د اړيکې $\pSat{v}{!A}$ پر بنسټ بيا د تاوتولوژۍ،
استلزام او د صدق وړتيا بنسټيز معنايي مفهومونه تعريفولے شو۔
!!^a{formula} هله تاوتولوژي، $\Entails !A$، ده چې هر !!{valuation}
ئې صادق کړي، يعنې د هر~$\pAssign{v}$ دپاره
$\pValue{v}(!A) = \True$ وي۔ د !!{formula}s يو سټ هله هغې ته
استلزام ورکوي، $\Gamma \Entails !A$، چې هر هغه !!{valuation} چې د
~$\Gamma$ ټول !!{formula}s صادق کړي، $!A$ هم صادق کړي۔ او د
!!{formula}s يو سټ هله د صدق وړ دے چې کوم !!{valuation} ئې ټول
!!{formula}s په يو وخت صادق کړي۔ ځکه چې !!{formula}s په استقرايي
ډول تعريف شوي او صدق بيا د !!{formula}s پر جوړښت د استقرا له مخې
تعريف شوے، موږ استقرا کارولے شو چې د خپلې معنٰی پوهنې خاصيتونه ثابت
او تعريف شوي معنايي مفهومونه يو له بل سره وتړو۔


\end{document}
